\documentclass[10pt,a4paper]{amsart}
\usepackage[margin=19mm]{geometry}
\usepackage{amsmath,amssymb,mathtools}
\usepackage[scr=rsfs,cal=euler]{mathalfa}
\usepackage{xfrac}
\usepackage[unicode,hidelinks,bookmarks=false]{hyperref}
\newcommand{\ie}{i.e.\ }
\newcommand{\cf}{cf.\ }
\newcommand{\resp}{respectively\ }
\newcommand{\Subsection}{\subsection}
\providecommand{\nobreakdash}{\nobreak\hbox{-}}
\setlength{\emergencystretch}{2em}
\multlinegap0pt
\usepackage{tikz}
\usetikzlibrary{cd}
\usepackage{amssymb}
\usepackage{dsfont}
\usepackage{mathtools}
\newtheorem{lemma}{Lemma}
\renewcommand{\H}{\textbf{H}}
\newcommand{\HBM}[2]{H_\bullet^{#1}(#2)}
\renewcommand{\Im}{\text{Im}}
\newcommand{\R}{\mathcal{R}}
\newcommand{\bbB}{\mathbb{B}}
\newcommand{\bbG}{\mathbb{G}}
\renewcommand{\subset}{\subseteq}
\newcommand{\tM}{\widetilde{M}}
\newcommand{\tW}{\mathbb{W}}
\title{Convolution algebras associated to representations}
\author{Dragoș Crișan}
\date{Source: arXiv:2606.19783v2 (CC BY 4.0)}
\begin{document}
\maketitle

\noindent\emph{Source and licence.} Convolution algebras associated to representations. Version arXiv:2606.19783v2 (CC BY 4.0), distributed by arXiv under the Creative Commons Attribution 4.0 International licence, http://creativecommons.org/licenses/by/4.0/, as recorded in the arXiv licence metadata for this exact version and shown on the abstract page https://arxiv.org/abs/2606.19783v2. Full licence notice: \texttt{licenses/arxiv-2606.19783-CC-BY-4.0.txt}.\par\smallskip

\noindent\textit{Editorial scope.} Counted unit: the article's lemma affine\_Image\_in\_HM (source0.tex lines 841--844). The article's complete original proof of it is delivered with it (source0.tex lines 846--859, one \texttt{proof} environment of 14 lines). No mathematics is written, repaired or re-derived: the statement and the proof are the article's own lines, in the article's own notation, and no retained line is substituted or re-flowed.  Retained uncounted as the source-local context the proof needs: lines 552--555; lines 817--824.  Nothing was omitted from any retained span, so the package carries no omission record.  No retained line contains a citation, so the wrapper carries no bibliography and no bibliography entry.  The retained lines contain the article's own diagram code, which the wrapper typesets with the standard tikz packages; no image file is delivered and the package declares no asset dependency.  Editorial formatting only, in addition: an AMS \texttt{amsart} preamble that restates the article's own declarations and macros used by the retained lines, namely the environments \texttt{lemma} on separate counters, together with the standard packages those lines need.  The retained environments print in this document as Lemma 1 in order of appearance, whereas the article prints the same items with the numbering of its own full document.  The complete native source of the article as distributed in the arXiv e-print for this exact version is bundled unchanged as \texttt{sources/203170/source0.tex}.
\begin{lemma}
\label{finite_Image_in_HM}
    $\Im(\theta_0) \subset \H_M$.
\end{lemma}

\begin{equation}
\label{affine_diagram_M_module}
\begin{tikzcd}[column sep=large, row sep=large]
\mathcal{R} \times \{0\} \arrow[d] & \bbG \times \{ 0 \} \arrow[l, "\tilde{p}_0"{above}] \arrow[r, "q_0"] \arrow[d] & \bbG \times^{\bbB} \{ 0 \} \arrow[r, "m_0"] \arrow[d] & \{ 0 \} \arrow[d] \\
\mathcal{R} \times M \arrow[d, "i \times \text{id}"{left}] & p_M^{-1} (\R \times M) \arrow[l, "\tilde{p}_M"{above}] \arrow[r, "\tilde{q}_M"] \arrow[d] & \R \arrow[r, "\tilde{m}_M"] \arrow[d, "i"] & M \arrow[d] \\
\tM \times M & \bbG \times M \arrow[l, "p_M"{above}] \arrow[r, "q_M"{above}] & \tM \arrow[r, "m_M"{above}] & V,
\end{tikzcd}
\end{equation}

\begin{lemma}
\label{affine_Image_in_HM}
    $\Im(\theta_0) \subset \H_M$.
\end{lemma}

\begin{proof}
    Let $y\geq w$ in $\tW$. As in the proof of Lemma \ref{finite_Image_in_HM}, it is enough to prove that the following diagram commutes:
    \[\begin{tikzcd}
	\HBM{\bbB}{\R^y_{\leq w}} \otimes \HBM{\bbB}{0} && \HBM{\bbB}{0} \\
	\\
	\HBM{\bbB}{\R^y_{\leq w}} \otimes \HBM{\bbB}{M^y} && \HBM{\bbB}{M^y}&& 
	\arrow["\star"', from=1-1, to=1-3]
	\arrow["\text{id} \otimes \varepsilon_*", from=1-1, to=3-1]
	\arrow["\varepsilon_*", from=1-3, to=3-3]
	\arrow["\star",from=3-1, to=3-3]
\end{tikzcd}\]

Here, $\varepsilon : \{ 0 \}\hookrightarrow M^y$ is the closed embedding of the origin inside the vector space $M^y$, and the horizontal maps are the action maps defined before. The commutativity of the diagram follows from proper base change along the Cartesian squares in the first two rows of diagram (\ref{affine_diagram_M_module}).
\end{proof}

\end{document}
